\PassOptionsToPackage{numbers}{natbib}
\documentclass[12pt,a4paper]{article}

% Page geometry (this is the key)
% \usepackage[
%   left=2.5cm,
%   right=2.5cm,
%   top=2.5cm,
%   bottom=2.5cm
% ]{geometry}
\usepackage[margin=2cm]{geometry}

% TikZ
\usepackage{tikz}
\usetikzlibrary{arrows.meta,shapes.misc,backgrounds,fit,positioning}

% Fonts (same as standard LaTeX / INRIA)
\usepackage[T1]{fontenc}
\usepackage{lmodern}

% Paragraph spacing similar to INRIA class
\setlength{\parskip}{0pt}
\setlength{\parindent}{15pt}

% Line spacing (default, explicit)
\linespread{1.0} % single-spaced per template

% \documentclass{article}
% % Layout
% \usepackage[a4paper,margin=1in,headsep=10pt]{geometry}
% % Encoding and typography
% \usepackage[T1]{fontenc}
\usepackage{newtxtext}   % Times text, as required by the ANRF template
\let\hex\undefined % newtxtext defines \hex; our macros.tex redefines it
\usepackage{sourcesanspro} % sans only for headings
\usepackage{microtype}
\usepackage{setspace}
% \setstretch{1.08}
% \usepackage[backend=biber,style=numeric,sorting=nyt,defernumbers=true]{biblatex}

\usepackage{natbib}
\usepackage{bibentry}
\usepackage{hanging}
\usepackage{xcolor}
\newcommand\publication[1]{%
    \smallskip\par\hangpara{1.5em}{1}\begin{NoHyper}\bibentry{#1}\end{NoHyper}\smallskip
}
\makeatletter
\renewcommand{\bibsection}{\section*{\refname}}
\makeatother

\definecolor{owncitecolor}{RGB}{178, 102, 255}
\newcommand{\owncite}[1]{\textcolor{linkblue}{\textbf{\cite{#1}}}}
% Define custom colors using HEX codes
\definecolor{linkblue}{HTML}{007BFF} % A nice, modern blue
\definecolor{linkred}{HTML}{C00000}  % A dark, professional red for internal links
\definecolor{linkgreen}{HTML}{008000} % A nice dark green for citations

% Option 2: Prevent amsmath from defining these commands
% \usepackage{etoolbox}
\csletcs{iint}{relax}
\csletcs{iiint}{relax}
\csletcs{iiiint}{relax}
\csletcs{idotsint}{relax}

% \usepackage{amsmath}
\usepackage{amsmath,amssymb}

\usepackage{array}
\usepackage{subcaption}
\usepackage{graphicx}
\usepackage{tikz}
\setlength{\parindent}{2em}
\setlength{\parskip}{0em}
\usepackage{footnote}
\usepackage{tablefootnote}
% \usepackage[table]{xcolor}
\usepackage{float}
\usepackage{fancybox}
% \usepackage{pseudocode}
\usepackage{algorithm}
\usepackage[noend]{algpseudocode}
\usepackage{pdflscape}
\usepackage{pdfpages}
\usepackage{afterpage}
\usepackage{capt-of}
\usepackage{adjustbox}
\usepackage{pagecolor}
\usepackage{booktabs}
\usepackage{pdfpages}
\usepackage{hyperref}
\usepackage[
n,
operators,
advantage,
sets,
adversary,
landau,
probability,
notions,    
logic,
ff,
mm,
primitives,
events,
complexity,
asymptotics,
keys]{cryptocode}

\usepackage{graphicx}
\usepackage{xspace}
\usepackage{collectbox}
\usepackage{multicol}
\usepackage{multirow}
\usepackage{setspace}
\setlength{\columnsep}{1cm}
\usepackage[capitalise]{cleveref}
\crefname{figure}{Figure}{Figures}
\crefname{equation}{Equation}{Equations}
\usepackage{makecell}

\definecolor{DeepBlue}{RGB}{0, 0, 139}
\definecolor{DeepGreen}{RGB}{0, 100, 0}
\definecolor{DeepRed}{RGB}{139, 0, 0}
\definecolor{DeepPurple}{RGB}{85, 26, 139}
\definecolor{DarkCyan}{RGB}{0, 139, 139}
\definecolor{DarkSlateGray}{RGB}{47, 79, 79}
\definecolor{SaddleBrown}{RGB}{139, 69, 19}
\definecolor{Indigo}{RGB}{75, 0, 130}
\definecolor{Maroon}{RGB}{128, 0, 0}
\definecolor{Olive}{RGB}{85, 107, 47}
\definecolor{DeepBurgundy}{RGB}{80, 0, 20}
\definecolor{InkBlue}{RGB}{0, 0, 90}
\definecolor{ForestShadow}{RGB}{0, 45, 0}
\definecolor{DeepCharcoal}{RGB}{40, 40, 40}
\definecolor{AcademicPurple}{RGB}{50, 0, 80}

% Define Ubuntu terminal colors
\definecolor{ubuntubg}{RGB}{48, 10, 36}  % Slight brownish background
\definecolor{ubuntuorange}{RGB}{193, 120, 36}  % Ubuntu terminal orange
\definecolor{terminaltext}{RGB}{220, 220, 220}  % Light gray text
\definecolor{darkgraytext}{RGB}{51, 51, 51}

% Themed boxes
\usepackage[most]{tcolorbox}
\NewTColorBox{genericbox}{m O{accent}}{
title=#1,
colback=white,
colframe=#2,
boxrule=1pt,
borderline west={3pt}{0pt}{#2},
fonttitle=\sffamily\bfseries
}

% Colored summary box: \begin{genericbox}{Title}[color]
% \NewDocumentEnvironment{genericbox}{m O{blue!50}}{%
%   \colorlet{genericboxframe}{#2}%
%   \colorlet{genericboxbg}{genericboxframe!15!white}%
%   \begin{tcolorbox}[
%     title={#1},
%     colframe=genericboxframe,
%     colback=genericboxbg,
%     fonttitle=\bfseries\sffamily,
%     sharp corners,
%     boxrule=0.6pt,
%     left=4pt, right=4pt, top=3pt, bottom=3pt
%   ]%
% }{%
%   \end{tcolorbox}%
% }


\newcommand{\acknowledgementstyle}[2][black]{\textcolor{#1}{\textit{#2}}}

\newtheorem{exmp}{Example}[section]
% \newtheorem{exmp1}{Example}[section]
% \newenvironment{exmp}
%   {\color{red} \begin{exmp1}}
%   {\end{exmp1}}

\newtheorem {lemma} {Lemma}
\newtheorem {proposition} {Proposition}
\newtheorem {theorem} {Theorem}
\newtheorem {corollary} {Corollary}
\newtheorem {definition} {Definition}
\newtheorem {remark} {Remark}

\newcommand{\ie}{{i.e.}\xspace}
\newcommand{\eg}{{e.g.}\xspace}

\newcommand{\malg}[1]{\textsf{Matsui’s Algorithm-#1}}

\usetikzlibrary{positioning}
%%% Custom TikZ addons
\usetikzlibrary{crypto.symbols}
%\tikzset{shadows=no}  
\usetikzlibrary{decorations.pathreplacing}
\usetikzlibrary{arrows}
\usetikzlibrary{arrows.meta}
\usetikzlibrary{patterns}
\usetikzlibrary{shapes}
\usetikzlibrary{shadows}
\usetikzlibrary{calc}
\usetikzlibrary{math}
\newcommand\blankpage{%
    \null
    \thispagestyle{empty}%
    \addtocounter{page}{-1}%
    \newpage}

% \newenvironment{proof}{\noindent{\bf Proof:}\newline\indent}%
%                       {\hfill $\Box$\par}
\newenvironment{miniproof}{\noindent{\bf Sketch Of Proof:}\newline\indent}%
                      {\hfill $\Box$\par}


\newcommand{\func}[1]{\operatorname{\mathit{#1}}}  % or \mathrm
\newcommand{\const}[1]{\mathit{#1}}                % or \mathrm
\newcommand{\wsub}[1]{\mathrm{#1}}
\newcommand{\bv}[1]{{\mathbf{#1}}}
\newcommand{\ppp}{\mathsf{p}}
\newcommand{\exr}{\oplus}
\newcommand{\etal}{\textit{et al.}}
\newcommand{\sbb}{\textsf{SBox}\xspace}
\newcommand{\tj}{\textsf{TinyJAMBU}\xspace}
\newcommand{\acpkm}{\ensuremath{\mathsf{ACPKM}}\space}
\newcommand{\ctracpkm}{\ensuremath{\mathsf{CTR}\mbox{-}\mathsf{ACPKM}}\space}
\newcommand{\h}[1]{\ensuremath{\mathsf{H}_{#1}}}
\newcommand{\saeb}{\textsf{SAEB}\space}
\newcommand{\parse}[1]{\ensuremath{\stackrel{{#1}}{\gets}}}
\newcommand{\bits}{\{0,1\}}
\newcommand{\tx}{\texttt}
\newcommand{\txf}{\textsf}
\newcommand{\cx}{\textsf}
\newcommand{\calO}{\mathcal{O}}
\newcommand{\pr}{\ensuremath{\textsf{Pr}}}

\definecolor{nvidiagreen}{RGB}{76, 175, 80}
\definecolor{wildstrawberry}{RGB}{255, 67, 164}
\definecolor{RoyalBlue}{RGB}{65,105,225}
\newcommand{\red}{\textcolor{wildstrawberry}}
\newcommand{\blue}{\textcolor{RoyalBlue}}
\newcommand{\olive}{\textcolor{olive}}
\newcommand{\orange}{\textcolor{orange}}
\newcommand{\violate}{\textcolor{violate}}

%Colors (used by Shibam)
\newcommand{\tcr}{\textcolor{red}}
\newcommand{\tcb}{\textcolor{blue}}
\newcommand{\tcg}{\textcolor{green}}

\newcommand{\done}{\textbf{\textcolor{red}{DONE}}\xspace}
\newcommand{\ntc}{\textbf{\textcolor{red}{NEED TO CHECK}}\xspace}
\newcommand{\ntd}{\textbf{\textcolor{red}{NEED TO DISCUSS}}\xspace}

\newcommand{\currentmonthyear}{\monthname[\month] \number\year}
% Macro for the current date in MM/DD/YYYY format
\newcommand{\currentdate}{%
  \ifnum\month<10 0\fi\number\month/% Add leading zero for month
  \ifnum\day<10 0\fi\number\day/% Add leading zero for day
  \number\year% Year
}

%%%%%%%%%%%%%%%%%%%%%%%%%%%%%%%%%%%MACROS FFT PAPER%%%%%%%%%%%%%%%%%%%%%%%
% recurrent symbols/name...
\newcommand{\GF}{\mathbb{F}_2}
\newcommand{\indic}{\boldone}
\newcommand{\emm}{\textsf{EM}\xspace}
\newcommand{\xx}{\mathsf{X}\xspace}

%\usepackage{bbm}
%\newcommand{\indic}{\mathbbm{1}}
% ciphers
\newcommand{\AES}{\texttt{AES}\xspace}
\newcommand{\sAES}{\texttt{small-AES}\xspace}
\newcommand{\MISTY}{\texttt{MISTY1}\xspace}

% aes subfunctions
\newcommand{\ak}{\textsc{AddRoundKey}\xspace}
\newcommand{\subB}{\textsc{SubBytes}\xspace}
\newcommand{\sr}{\textsc{ShiftRows}\xspace}
\newcommand{\mc}{\textsc{MixColumns}\xspace}

%Hexadecimal constants (field elements)
\newcommand{\hex}[1]{\ensuremath{\mathtt{#1}_{\mathsf{x}}}}

\newcommand{\pricePartial}{497\xspace} %in dollar
\newcommand{\priceFFT}{418\xspace} %in dollar
\newcommand{\priceOur}{5\xspace} %in dollar
\newcommand{\timePartial}{4859\xspace} %in minutes
\newcommand{\timeFFT}{3120\xspace} %in minutes
\newcommand{\timeOur}{48\xspace} %in minutes
\newcommand{\speedup}{65\xspace} % minimal speed-up in real time between our method compared to FFT and original partial sum
%\newcommand{\cheapup}{62.8\xspace} % minimal speed-
\newcommand{\cheapup}{83\xspace} % minimal speed-


%%%%%%%%%%%%%%%%%%%%%%%%%%%%%%%%%%%MACROS Tinyjambu PAPER%%%%%%%%%%%%%%%%%%%%%%%
\newcommand{\Pf}[1][1024]{\ensuremath{ \mathcal{P}^{#1}}\xspace}
%%%%%%%%%%%%%%%%%%%%%%%%%%%%%%%%%%%MACROS RUP PAPER%%%%%%%%%%%%%%%%%%%%%%%
\newcommand{\saeaes}{\textsf{SAEAES}\xspace}
\newcommand{\rv}[1]{\ensuremath{\mathsf{#1}}}
\newcommand{\bbN}{\mathbb{N}}
\newcommand{\ozs}{\textsf{ozs}\xspace}
\newcommand{\ver}{\ensuremath{\mathsf{ver}}}
\newcommand{\advA}{\mathscr{A}}
\newcommand{\ac}{\ensuremath{\mathsf{\Theta}}}
\newcommand{\intrupadvt}{\ensuremath{\advt^{\mathsf{int\text{-}rup}}}}
\newcommand{\advt}{\ensuremath{\mathbf{Adv}}}
\renewcommand{\enc}{\ensuremath{\mathsf{enc}}}
\renewcommand{\dec}{\ensuremath{\mathsf{dec}}}
\newcommand{\abort}{\ensuremath{\rv{\bot}}}
\newcommand{\rerv}{\ensuremath{\rv{\Lambda}_1}}
\newcommand{\idrv}{\ensuremath{\rv{\Lambda}_0}}
\renewcommand{\event}[1]{\ensuremath{\mathtt{#1}}}
\renewcommand{\bad}{\ensuremath{\event{bad}}}
\newcommand{\good}{\ensuremath{\event{good}}}
\newcommand{\A}{\mathcal{A}}
\newcommand{\pf}{\noindent{\bf Proof. }}
\newcommand{\defn}{\stackrel{\mathrm{\Delta}}{=}}
\newcommand{\mx}{\mathsf}
\newcommand{\E}{\mathbf{E}}
\newcommand{\rvv}[1]{\ensuremath{\rv{\bar{#1}}}}
\newcommand{\init}{\ensuremath{\mathsf{Init}}\xspace}

%%%%%%%%%%%%%%%%%%%%%%%%%%%%From QARMA-PAC%%%%%%%%%%%%%%%%%%%%%%%%%%%%%%
\newcommand{\deltain}{\delta_{\mathit{in}}}
%\newcommand{\deltaT}{\delta_{\mathrm{T}}}
\newcommand{\deltaout}{\delta_{\kern-.055em\mathit{out}}}
\newcommand{\Deltain}{\Delta_{\mathit{in}}}
\newcommand{\Deltaout}{\Delta_{\kern-.027em\mathit{out}}}
\newcommand{\Nablain}{\nabla_{\mathit{in}}}
\newcommand{\Nablaout}{\nabla_{\kern-.027em\mathit{out}}}
\newcommand{\DeltaT}{\Delta_{\mathrm{T}}}
\newcommand{\din}{d_{\mathit{in}}}
\newcommand{\dinstar}{\din^{\kern.04em*}}
\newcommand{\dout}{d_{\kern-.027em\mathit{out}}}
\newcommand{\dT}{d_{\kern-.027em\mathit{T}}}
\newcommand{\hatdout}{{\widehat d}_{\kern-.027em\mathit{out}}}
\newcommand{\Din}{D_{\mathit{in}}}
\newcommand{\Dout}{D_{\kern-.11em\mathit{out}}}
\newcommand{\DT}{D_{\mathrm{T}}}
\newcommand{\hatDout}{{\widehat D}_{\kern-.11em\mathit{out}}}

\newcommand{\smin}{s_{\mathit{min}}}
\newcommand{\zetamin}{\zeta_{\mathit{min}}}
\let\mycal\mathcal

\newcommand{\calK}{\mycal{K}}
\newcommand{\rin}{r_{\mathit{in}}}
\newcommand{\rout}{r_{\kern-.11em\mathit{out}}}
\newcommand{\rmid}{r_{\mathit{mid}}}
\newcommand{\Kbits}{\calK}
\newcommand{\Kin}{\Kbits_{\mathit{in}}}
\newcommand{\Kout}{\Kbits_{\kern-.082em\mathit{out}}}
\newcommand{\keyguess}{k}
\newcommand{\kin}{\keyguess_{\mathit{in}}}
\newcommand{\kout}{\keyguess_{\kern-.055em\mathit{out}}}
\newcommand{\cin}{c_{\mathit{in}}}
\newcommand{\cout}{c_{\kern-.055em\mathit{out}}}

\newcommand{\TGF}{T_{\textit{GF}}}
\newcommand{\SGF}{S_{\textit{GF}}}
\newcommand{\TMR}{T_{\textit{MR}}}
\newcommand{\TMW}{T_{\textit{MW}}}
\newcommand{\Tenc}{T_{\calE}}
\newcommand{\Tcollect}{T_{\textit{C}}}
\newtheorem{example}{Example}

%%%%%%%%%%%%%%%%%%%%%%%%%%%%%%ZKPS%%%%%%%%%%%%%%%%%%%%%%%%%%%%%%
\newcommand{\calP}{\mathcal{P}}
\newcommand{\calV}{\mathcal{V}}

\newcommand{\FA}{\textsf{FA}\xspace}
\newcommand{\FAs}{\textsf{FA}s\xspace}
% \newcommand{\XOR}{\textsf{XOR}\xspace}
\newcommand{\XORs}{\textsf{XOR}s\xspace}



% Colors
% \usepackage[dvipsnames]{xcolor}
% Theme palette (tweak freely)
% \definecolor{theme}{HTML}{0E5A8A}   % primary
\definecolor{theme}{HTML}{0E5A8A}   % primary
\definecolor{accent}{HTML}{F39C12}  % accent
\definecolor{softbg}{HTML}{F7FBFF}  % subtle background
\definecolor{linkcol}{HTML}{0A6AA1} % URL/link color
\definecolor{citecol}{HTML}{0B8A6B} % citations
% Lists
\usepackage[shortlabels]{enumitem}
% \setlist{nosep}
% \renewcommand\labelitemi{\textcolor{theme}{$\blacktriangleright$}}

% Headers/footers
\usepackage{fancyhdr}
\usepackage{lastpage}
\pagestyle{fancy}
\fancyhf{}

% Colored thin head rule
\makeatletter
\renewcommand{\headrule}{\hbox to\headwidth{\color{theme!50}\leaders\hrule height 0.4pt\hfill}}
\makeatother
\renewcommand{\headrulewidth}{0.4pt}
\renewcommand{\footrulewidth}{0pt}

% Section titles
\usepackage[explicit]{titlesec}
\titleformat{\section}
{\Large\sffamily\bfseries\color{theme}}
{\thesection}{0.75em}{#1}
[\color{theme!50}\titlerule]
\titleformat{\subsection}
{\large\sffamily\bfseries\color{theme!85!black}}
{\thesubsection}{0.6em}{#1}
\titleformat{\subsubsection}
{\normalsize\sffamily\bfseries\color{theme!70!black}}
{\thesubsubsection}{0.5em}{#1}
\titlespacing*{\section}{0pt}{1.4ex plus .3ex}{0.9ex}
\titlespacing*{\subsection}{0pt}{1.0ex plus .2ex}{0.5ex}
\titlespacing*{\subsubsection}{0pt}{0.8ex}{0.4ex}

\newcommand{\yourname}{Shibam Ghosh}
\newcommand{\youremail}{shibam@cse.iith.ac.in}
\newcommand{\yourweb}{shibamcrs.github.io}
% Hyperlinks (load last)
\hypersetup{
unicode=true,
colorlinks=true,
linkcolor=theme,
citecolor=citecol,
urlcolor=linkcol,
pdfstartview=FitV,
pdftitle={\yourname},
pdfauthor={\yourname}
}
% Header content (after hyperref so links work)
% \fancyhead[C]{\footnotesize\sffamily\bfseries
% \yourname\quad
% \textcolor{linkcol}{\href{https://\yourweb}{\yourweb}}\quad
% \textcolor{linkcol}{\href{mailto:\youremail}{\youremail}}
% }
\fancyfoot[C]{\footnotesize\sffamily Page \thepage\ of \pageref{LastPage}}

\newcommand{\soptitle}{DARC: Design and Analysis of Rapid Cryptography}
\newcommand{\sopsubtitle}{Full Proposal -- Technical Document, Category II (Use-Inspired Basic Research)}
\begin{document}
\begin{center}
\vspace*{-0.8cm}
{\sffamily\bfseries\Large \textcolor{theme}{\soptitle}}\par
\vspace{2pt}
{\sffamily\sopsubtitle}\\[4pt]
Shibam Ghosh\textsuperscript{1} (PI) \quad Dhiman Saha\textsuperscript{2} (Co-PI) \\[1ex]
\small \textsuperscript{1}Indian Institute of Technology Hyderabad
  \small \textsuperscript{2}Indian Institute of Technology Bhilai \\[1ex]
% {\small \yourname, Dhiman Saha}
\vspace{8pt}
\hrule
\end{center}

% ─────────────────────────────────────────────────────────────────────────────
\section{Use-Inspired Context, Scientific Rationale \& Knowledge Gap}
% ─────────────────────────────────────────────────────────────────────────────

The design of symmetric-key cryptographic primitives has historically been heavily influenced by the specific requirements of their target applications.
A cryptographic primitive, such as a block cipher, hash function, or pseudorandom function, that is optimized for one context may be ill-suited for others, making the process a balancing act between security and efficiency dictated by each target environment.
Thus, different applications, ranging from high-throughput primitives for network communications to lightweight ciphers for constrained devices, have led to different corners of the design space being explored.
This diversity in requirements underscores the importance of designing new symmetric-key primitives that are well-adapted to their intended applications.
This proposal focuses on one such domain with unique requirements: symmetric-key primitives for applications demanding \emph{low-latency} cryptographic operations.
Figure~\ref{fig:design-space} situates this domain within the broader landscape of special-purpose symmetric-key cryptography.

\begin{figure}[htbp]
\centering
\begin{tikzpicture}[scale=0.72, font=\sffamily]

\tikzset{
  rtitle/.style  = {font=\small\sffamily\bfseries, align=center},
  rciphers/.style= {font=\fontsize{6.2}{7.5}\selectfont\sffamily\itshape, align=center},
  rapplic/.style = {font=\fontsize{5.8}{7}\selectfont\sffamily, text=gray!70, align=center},
  rnote/.style   = {font=\fontsize{6}{7}\selectfont\sffamily\itshape, align=center},
}

%% ── AXES ────────────────────────────────────────────────────────────────────
\draw[-{Stealth[length=6pt]}, line width=0.9pt] (0,0) -- (14,0)
  node[right, font=\small\sffamily\bfseries] {Hardware Area / Resources};
\draw[-{Stealth[length=6pt]}, line width=0.9pt] (0,0) -- (0,11.2)
  node[above, font=\small\sffamily\bfseries] {Latency};

% Axis endpoint labels
\node[font=\fontsize{6}{7}\selectfont\sffamily, gray, below=2pt] at (1.3, 0)  {Constrained};
\node[font=\fontsize{6}{7}\selectfont\sffamily, gray, below=2pt] at (12.5, 0) {Unconstrained};
\node[font=\fontsize{6}{7}\selectfont\sffamily, gray, rotate=90, left=8pt]  at (0, 1.2) {$\approx\!1$ cycle};
\node[font=\fontsize{6}{7}\selectfont\sffamily, gray, rotate=90, left=8pt]  at (0, 9.8) {Many cycles};

% 1-cycle dashed reference
\draw[gray!25, dashed, line width=0.8pt] (0.2, 1.6) -- (13.5, 1.6);
\node[font=\fontsize{5.2}{6}\selectfont\sffamily, gray!55, anchor=east] at (13.4, 1.78) {1-clock-cycle target};

%% ── LIGHTWEIGHT CRYPTO ──────────────────────────────────────────────────────
\fill[green!12, draw=green!50!black, line width=1.1pt, rounded corners=11pt]
  (0.3, 3.8) rectangle (4.3, 10.9);
\node[rtitle, text=green!38!black] at (2.3, 10.5) {Lightweight};
\node[rciphers, text=green!32!black] at (2.3, 9.1)
  {GIFT, PRESENT\\SIMON, SPECK\\ASCON, PHOTON\\SKINNY};
\node[rapplic] at (2.3, 7.4)
  {\emph{IoT, RFID,}\\[1pt]\emph{smart cards,}\\[1pt]\emph{sensor networks}};
\node[rnote, text=green!38!black] at (2.3, 4.3)
  {Small area,\\tolerates latency};

%% ── GENERAL-PURPOSE ─────────────────────────────────────────────────────────
\fill[orange!11, draw=orange!58!black, line width=1.1pt, rounded corners=11pt]
  (4.6, 3.8) rectangle (9.0, 8.7);
\node[rtitle, text=orange!55!black] at (6.8, 8.3) {General-Purpose};
\node[rciphers, text=orange!43!black] at (6.8, 7.1)
  {AES, SHA-3\\ChaCha20, 3DES};
\node[rapplic] at (6.8, 5.7)
  {\emph{TLS, cloud services,}\\[1pt]\emph{storage, general}\\[1pt]\emph{computing}};
\node[rnote, text=orange!50!black] at (6.8, 4.3) {Balanced trade-offs};

%% ── HIGH-THROUGHPUT (PIPELINED) ─────────────────────────────────────────────
\fill[violet!10, draw=violet!52!black, line width=1.1pt, rounded corners=11pt]
  (9.3, 3.8) rectangle (13.5, 8.7);
\node[rtitle, text=violet!55!black] at (11.4, 8.3) {High-Throughput};
\node[rciphers, text=violet!42!black] at (11.4, 7.1)
  {AES-NI, AES-GCM\\Keccak (HW)\\pipelined AES};
\node[rapplic] at (11.4, 5.7)
  {\emph{Network, VPN,}\\[1pt]\emph{disk encryption,}\\[1pt]\emph{bulk data}};
\node[rnote, text=violet!48!black] at (11.4, 4.3)
  {Pipelining $\not\!\Rightarrow$ low latency};

%% ── HE / ZK-FRIENDLY ────────────────────────────────────────────────────────
\fill[red!8, draw=red!48!black, line width=1.1pt, rounded corners=11pt]
  (4.6, 9.0) rectangle (13.5, 10.9);
\node[rtitle, text=red!52!black] at (9.0, 10.55) {HE\,/\,ZK-Friendly};
\node[rciphers, text=red!38!black] at (9.0, 9.8)
  {Poseidon,\ MiMC,\ Rescue-Prime,\ Griffin,\ Reinforced Concrete};
\node[rapplic] at (9.0, 9.3)
  {\emph{Homomorphic encryption, zero-knowledge proofs, blockchain}};

%% ── LOW-LATENCY CRYPTO (FOCUS) ──────────────────────────────────────────────
% outer glow
\fill[blue!6, rounded corners=15pt]  (9.2, 0.15) rectangle (13.6, 3.65);
% main box
\fill[blue!11, draw=theme, line width=2pt, rounded corners=12pt]
  (9.3, 0.25) rectangle (13.5, 3.55);

\node[rtitle, text=theme] at (11.4, 3.2) {Low-Latency Crypto};
\node[rciphers, text=blue!58!black] at (11.4, 2.3)
  {PRINCE, PRINCEv2, MANTIS\\QARMA, QARMAv2, SCARF\\ARADI, ChiLow};
\node[rapplic] at (11.4, 1.35)
  {\emph{Memory encryption, cache}\\[1pt]
   \emph{randomization, pointer auth (PAC),}\\[1pt]
   \emph{automotive, industrial control}};

% "This proposal" badge
\node[
  draw=theme, fill=theme!15, line width=0.9pt, rounded corners=4pt,
  font=\fontsize{6}{7}\selectfont\sffamily\bfseries, text=theme,
  inner sep=3pt
] at (11.4, 0.52) {$\bigstar$~~This Proposal};

%% ── EMPTY CORNER ANNOTATION (bottom-left: impossible region) ────────────────
\node[
  font=\fontsize{6.5}{8}\selectfont\sffamily, text=gray!40,
  align=center
] at (2.3, 2.7) {\emph{(ultra-small area}\\[-1pt]\emph{+ ultra-low latency}\\[-1pt]\emph{is infeasible)}};

%% ── CONNECTOR: low-latency to high-throughput (latency wall) ────────────────
\draw[gray!40, dashed, {Stealth[length=4pt]}-{Stealth[length=4pt]}, line width=0.7pt]
  (11.4, 3.6) -- (11.4, 3.8);
\node[font=\fontsize{5.5}{6}\selectfont\sffamily, gray!55] at (11.4, 3.72) {};

\end{tikzpicture}
\caption{Design space of special-purpose symmetric-key cryptographic primitives.
  The horizontal axis measures hardware resource cost (gate area); the vertical axis
  measures per-operation latency.
  Low-latency crypto occupies the bottom-right corner: it accepts a large hardware
  footprint (fully-unrolled circuit) in exchange for near-single-cycle latency,
  a property that high-throughput pipelined designs cannot provide.}
\label{fig:design-space}
\end{figure}

% \subsection{Primitives for Low-Latency Crypto}
\section{Research Objectives, Hypotheses \& Success Metrics}

Latency, or the time delay from input to output, is critical especially in applications requiring hardware-assisted mechanisms to protect against software exploitation.
Memory and code protection stand out as crucial: while in contemporary processor architectures it is common to protect the CPU within secure enclaves to safeguard sensitive computations, extending such protections to any external memory remains impractical due to physical and architectural constraints.
As a result, \emph{data-in-use} in RAM remains a common target for physical attacks such as the \emph{Cold Boot} attack~\cite{DBLP:conf/uss/HaldermanSHCPCFAF08}, and code stored in external memory is vulnerable to eavesdropping.
This underscores the importance of schemes that protect cryptographic memory and code.
To be practical in high-performance systems, the primitives used in this context must offer minimal added latency.

Other latency-sensitive applications include cache randomization (to protect against cache side-channel attacks, e.g., \textsf{PRIME+PROBE}~\cite{DBLP:journals/iacr/OsvikST05} and \textsf{FLUSH+RELOAD}~\cite{DBLP:conf/uss/YaromF14}), pointer authentication (\tx{PAC} feature~\cite{ARMv8-A-2016-additions} in Arm~v8), automotive communication, and industrial control networks.
As secure processor architectures continue to evolve, we expect such features to become increasingly prevalent, particularly as more highly optimised cryptographic primitives become available.
While throughput can benefit from pipelining and parallelism, designing a secure encryption algorithm whose hardware implementation is fast remains fundamentally challenging.
In a 1-cycle architecture, what matters is the latency of a fully-unrolled hardware implementation --- one where each round of a block cipher is mapped onto distinct and simultaneous hardware logic and the entire computation executes within a single clock cycle.
Although the recently completed NIST lightweight encryption/authentication standardisation effort addressed small area and low-power budgets, the design of low-latency cryptographic primitives remains a relatively new and rapidly evolving research field.

\begin{genericbox}{Why This Problem Matters}[teal!60]
Low-latency symmetric-key primitives underpin a growing class of hardware security mechanisms, from transparent memory encryption and pointer authentication to secure caches and automotive real-time networks.
Because these mechanisms sit on the critical path of processor execution, any added latency directly degrades system performance.
Yet the security of the small, fast designs that meet this constraint is far less understood than that of conventional block ciphers.
This proposal addresses that gap directly.
\end{genericbox}

\subsection{Research Objectives and Hypotheses}

\paragraph{Objective 1: Automated, accelerated and hybrid cryptanalysis of low-latency ciphers.}
Develop CP/SAT/MILP-based methods that jointly search for distinguishers and key-recovery attacks on low-latency block ciphers.
Tailor the resulting attack algorithms for high-performance GPU implementation, so that solver-found attacks can be verified and pushed further in practice.
Combine machine-learning-based distinguishers with classical techniques in principled hybrid approaches.

\textbf{Hypothesis H1:}
Automated and hybrid cryptanalytic methods, backed by GPU-accelerated implementations, can reveal new weaknesses in existing low-latency ciphers or improve their best known analysis.

\textbf{Success metrics:}
(i) New distinguishers or key-recovery results on at least two of \tx{ARADI}, \tx{ChiLow}, \tx{SCARF} and \tx{QARMAv2}, each improving on the best published attack in rounds covered or in complexity.
(ii) At least one attack implemented on GPU, with a measured speed-up over its CPU reference.
(iii) A hybrid ML--classical distinguisher evaluated on at least one target cipher.
(iv) An open-source constraint-model library covering all four target ciphers.

\paragraph{Objective 2: Physical attacks on fully unrolled
implementations.}
Develop automated methods to identify fault-related distinguishers
and validate potential attacks experimentally.

\textbf{Hypothesis H2:}
Automated fault analysis can identify exploitable properties
of fully unrolled implementations and guide practical fault attacks.

\textbf{Success metrics:}
(i) A publicly released fault-analysis tool.
(ii) Experimental validation of at least one fault attack on hardware.
(iii) Evidence-based hardening guidelines.

\paragraph{Objective 3: Microarchitectural security.}
Analyse the cryptographic components of CPU memory-encryption
or pointer-authentication mechanisms under realistic threat models.

\textbf{Hypothesis H3:}
Microarchitectural threat modelling can reveal security limitations
not adequately captured by primitive-level analysis alone.

\textbf{Success metrics:}
(i) Security evaluation of at least one target mechanism.
(ii) Where warranted, cryptographic improvements or hardening recommendations supported by appropriate evidence.

\paragraph{Objective 4: Design criteria and an indigenous low-latency cipher.}
Formalise how cryptographic security relates to implementation characteristics (gate depth, area, fault and probing exposure of the unrolled circuit) and to the latency requirements of the target applications.
Use these criteria, together with the findings of Objectives~1--3, to design and evaluate \tx{DARC}: an indigenous low-latency block cipher tailored to secure-hardware and microarchitectural use.

\textbf{Hypothesis H4:}
Design criteria that jointly account for algorithmic security, physical-attack exposure and microarchitectural context lead to a low-latency primitive with a stronger security argument than existing designs, at comparable latency.

\textbf{Success metrics:}
(i) A documented set of design criteria, with implementation metrics compared against \tx{PRINCEv2}, \tx{QARMAv2} and \tx{ARADI}.
(ii) A complete public specification of \tx{DARC}, with reference software and unrolled hardware implementations.
(iii) A security analysis of \tx{DARC} covering the attack classes of Objectives~1--3, with latency and area benchmarked against the above baselines.

\paragraph{Objective 5: Reproducibility and research training.}
Develop reusable research artefacts and strengthen expertise
in symmetric-key cryptography, automated analysis, and hardware security.

\textbf{Hypothesis H5:}
Reproducible tools and documented experiments facilitate
independent verification and further research.

\textbf{Success metrics:}
(i) Open-source tools and datasets accompanying every publication, with documented experimental procedures.
(ii) Two PhD students trained, one at each institute: one on automated and ML-assisted cryptanalysis, one on fault analysis of unrolled implementations.
(iii) Two postdoctoral researchers trained, one at each institute: one driving the constraint-model and GPU toolchain, one driving the hardware experimentation and the \tx{DARC} implementation.
% ─────────────────────────────────────────────────────────────────────────────
\section{Scientific Novelty, Advancement of Knowledge \& Utility}
% ─────────────────────────────────────────────────────────────────────────────

The introduction of application-specific optimisations raises fresh challenges in security evaluation.
\emph{Cryptanalysis} plays a vital role in testing the strength of such cryptographic algorithms and in identifying weaknesses that adversaries might exploit.
Unlike public-key cryptosystems, whose security relies mainly on the presumed hardness of mathematical problems, symmetric-key schemes are evaluated by their resilience against a diverse and evolving set of attacks.
It is thus essential for designers to present clear security arguments and invite public scrutiny to establish and maintain \emph{trust} in their algorithms.

\subsection{State of the Art and Open Problems}

Recent years have seen the development of several block ciphers explicitly designed for low-latency applications:
\tx{PRINCE}~\cite{DBLP:conf/asiacrypt/BorghoffCGKKKLNPRRTY12},
\tx{PRINCEv2}~\cite{DBLP:conf/sacrypt/BozilovEKLLMNRT20},
\tx{MANTIS}~\cite{DBLP:conf/crypto/BeierleJKL0PSSS16},
\tx{QARMA}~\cite{DBLP:journals/tosc/Avanzi17},
\tx{QARMAv2}~\cite{DBLP:journals/tosc/AvanziBDEGNR23},
\tx{ChiLow}~\cite{DBLP:conf/eurocrypt/BelkheyarDGLMPRSSAV25},
\tx{SCARF}~\cite{DBLP:conf/uss/CanaleGLTTU23}, and
\tx{ARADI}~\cite{cryptoeprint:2024/1240}.
To achieve low latency, these designs deliberately adopt a small gate depth, a reduced number of rounds, and simpler key schedules compared to conventional cryptographic algorithms.
While this makes them efficient, it also concentrates the security burden on fewer components, making systematic cryptanalytic evaluation indispensable.

\subsection{Novelty and Contributions Beyond the Baseline}

\paragraph{Automated distinguisher and key-recovery search.}
Constraint-solving tools, such as CP, SAT, and MILP have proven powerful for uncovering new distinguishers and key-recovery attacks in conventional ciphers, but their systematic application to the low-latency cipher family is nascent.
Our prior work attacking \tx{QARMA}~\owncite{DBLP:journals/tosc/AvanziDG25} using SAT-based methods demonstrates the viability of this approach and establishes a concrete baseline.
A key novelty of this proposal is to jointly optimise distinguisher quality and key-recovery complexity, going beyond the typical practice of treating them as separate problems.

\paragraph{Machine-learning and classical hybrid cryptanalysis.}
Neural differential distinguishers have been explored for AES-like designs, yet their principled combination with classical techniques such as meet-in-the-middle and boomerang attacks remains largely open.
This proposal pursues a systematic investigation of such hybrid approaches, which have the potential to provide new insights into the security of low-latency cipher designs.

\paragraph{Physical attacks on fully-unrolled implementations.}
Fully unrolled hardware implementations present a threat model that is qualitatively different from that of iterative designs.
Clock-glitch fault injection is less effective because all rounds execute simultaneously, yet advanced techniques such as electromagnetic fault injection (EMFI) can still target specific spatial regions of the chip.
Furthermore, because all rounds are simultaneously present in silicon, the physical attack surface is enlarged, creating new opportunities for probing.
Our recent work on fault-invariant distinguishers~\owncite{Kundu_Ghosh_Aikata_Saha_2025} is, to our knowledge, the first to formalise this class of attacks; this proposal extends it by leveraging SAT/CP solvers to automate the search for fault-attack-friendly distinguishers and to improve post-processing key recovery from harvested hardware data.

\paragraph{Microarchitectural threat analysis.}
Hardware-based mitigations for important microarchitectural attacks such as memory encryption, pointer authentication, secure caches, increasingly rely on low-latency primitives.
We propose to systematically analyse the cryptographic components of these mitigation mechanisms to identify vulnerabilities specific to their microarchitectural deployment context.
This direction is novel in that it bridges the gap between the cryptographic and systems-security communities.

% ─────────────────────────────────────────────────────────────────────────────
\section{Research Approach, Methodology, Integration \& Feasibility}
% ─────────────────────────────────────────────────────────────────────────────

The research is structured into three temporally ordered thrusts, progressing from problems where established techniques provide a clear starting point (short-term) to problems that demand fundamentally new conceptual frameworks (long-term).
Feasibility at each stage is grounded in the team's existing expertise and preliminary results.

\begin{genericbox}{Research Summary}[teal!60]
This project is dedicated to the design and analysis of symmetric-key primitives optimised for applications requiring low-latency cryptographic operations.
Our expertise in symmetric-key design and cryptanalysis is well suited to investigating these challenges, systematically balancing security with application-specific performance.
Through this research, we aim to advance both the theoretical foundations and practical deployment of tailored primitives, bridging the gap between robust security guarantees and efficiency.
The resulting analysis and frameworks are intended to directly support the standardisation community by providing rigorous security analysis and clear metrics.
\end{genericbox}

\subsection{Thrust~1 (Short-term): Cryptanalysis of Existing Low-Latency Ciphers}

The immediate objective is to perform rigorous cryptanalytic evaluation of the current generation of low-latency block ciphers.
Automatic tools such as CP, SAT, and MILP are particularly powerful for uncovering new distinguishers and key-recovery attacks; we have used them extensively in prior work.
Leveraging such tools to identify the best distinguishers that facilitate optimal key-recovery attacks is a central methodological pillar.
Our SAT-based attack on \tx{QARMA}~\owncite{DBLP:journals/tosc/AvanziDG25} establishes both the feasibility and the productivity of this approach; we see significant opportunities to extend it to \tx{ARADI}~\cite{cryptoeprint:2024/1240}, \tx{ChiLow}~\cite{DBLP:conf/eurocrypt/BelkheyarDGLMPRSSAV25}, \tx{SCARF}~\cite{DBLP:conf/uss/CanaleGLTTU23}, and related designs.

In parallel, we will investigate the use of machine learning for the analysis of cryptographic algorithms.
While not entirely new as a direction, many open questions remain, particularly regarding the combination of ML-based distinguishers with classical cryptanalytic techniques in principled hybrid approaches.
Such combinations have the potential to complement traditional tools and provide new security insights into modern low-latency cipher designs.

Solver-found attacks are only as credible as their verification.
We will therefore tailor the key-recovery procedures to high-performance GPU implementation, exploiting the bit-sliced structure and shallow depth of low-latency ciphers.
This serves two purposes: experimentally confirming the complexities predicted by the models, and pushing attacks beyond what CPU implementations can reach in practice.
The GPU workstations already available at IIT~Bhilai (Section~7.4) are sufficient for this task.

\subsection{Thrust~2 (Medium-term): Physical Attacks on Low-Latency Implementations}

Physical security threats emerge at the circuit level, where attackers can directly observe or manipulate computations.
Unprotected implementations of low-latency cryptographic primitives are particularly vulnerable to side-channel exploitation.
Integrating countermeasures such as masking into fully unrolled hardware is challenging and often incurs considerable performance or area overhead.

Although fully unrolled implementations are generally resistant to clock-glitch fault injection (since all rounds execute in a single cycle), advanced techniques such as EMFI still allow attackers to target specific chip regions.
Moreover, having all rounds simultaneously present in hardware increases the physical ``surface area'', expanding opportunities for probing.
Deploying comprehensive protections such as duplication or physical shielding across the entire unrolled circuit is complex and costly.

We therefore aim to identify the most vulnerable components within these implementations.
Building on our prior work introducing fault-invariant distinguishers~\owncite{Kundu_Ghosh_Aikata_Saha_2025}, we will apply SAT/CP/MILP solvers to automate the discovery of fault-attack-friendly distinguishers and to perform key recovery from actual hardware measurements during post-processing.
This approach could significantly advance the ability to analyse and harden low-latency primitives against fault attacks.

\subsection{Thrust~3 (Long-term): Microarchitectural Security of CPU Cryptographic Features}

Security of modern CPUs has received significant attention in the wake of diverse microarchitectural attacks.
Hardware-based mitigations increasingly rely on encrypted communication between the CPU and its surrounding components, using techniques such as memory encryption, pointer authentication, and secure caches, all of which require low-latency primitives.
The long-term objective is to systematically analyse the cryptographic sub-systems of various CPU architectures, identify potential security vulnerabilities arising from their microarchitectural context, and develop efficient cryptographic solutions that safeguard against emerging threats.
Concrete initial targets are Arm pointer authentication (\textsf{FEAT\_PACQARMA3}, which the PI has already analysed~\owncite{DBLP:journals/tosc/AvanziDG25}), memory-encryption engines of the kind \tx{ARADI} was designed for, and cache-randomisation schemes built on \tx{SCARF}.
For each, we will state the threat model that the deployment context actually implies (truncated outputs, tweak and address reuse, attacker-controlled plaintexts, fault and probing access) and evaluate the primitive under that model rather than in isolation.

\subsection{Thrust~4 (Long-term): Design Criteria and the \tx{DARC} Cipher}

The analysis thrusts feed a design thrust.
Thrust~1 tells us which structural choices (S-box depth, linear-layer branch number, key-schedule simplicity) leave the most room for attacks.
Thrust~2 tells us which parts of an unrolled circuit are most exposed to faults and probing.
Thrust~3 tells us what the deployment context demands of the primitive.
We will consolidate these findings into explicit design criteria and quantify the security--latency trade-off they imply, benchmarked against \tx{PRINCEv2}, \tx{QARMAv2} and \tx{ARADI}.
Using these criteria we will design \tx{DARC}, an indigenous low-latency block cipher.
The design will be accompanied by self-analysis with the tools of Thrust~1, by unrolled hardware implementations and latency/area figures from IIT~Bhilai, and by a public specification inviting third-party cryptanalysis.
This thrust starts only in the second half of the project, once the analysis results are available.

\subsection{Work Packages, Integration and Risk Management}

Table~\ref{tab:wp} maps the thrusts onto six work packages, each with a lead institute, the objective it serves, and its duration.
Every technical work package has a clear hand-over: WP1 and WP2 deliver attack models and tools that WP3 reuses for fault-friendly distinguishers; WP3 and WP4 deliver the exposure and threat-model data that WP5 consumes.

\begin{table}[htbp]
\centering\small
\begin{tabular}{@{}p{0.9cm}p{5.5cm}p{1.8cm}p{1.1cm}p{1.4cm}p{3.6cm}@{}}
\toprule
\textbf{WP} & \textbf{Scope} & \textbf{Lead} & \textbf{Obj.} & \textbf{Months} & \textbf{Main outputs} \\
\midrule
WP1 & Constraint models; joint distinguisher and key-recovery search; GPU-accelerated attack implementations & IITH & O1 & 1--30 & D1, D2 \\
WP2 & Neural distinguishers; ML--classical hybrid framework & IITH & O1 & 7--36 & D3 \\
WP3 & Fault-friendly distinguisher search tool; unrolled hardware targets; EMFI and ChipWhisperer campaigns; hardening guidelines & IIT Bhilai & O2 & 7--42 & D4, D5, D6 \\
WP4 & Threat modelling and analysis of PAC, memory-encryption and cache-randomisation mechanisms; countermeasures & joint & O3 & 19--48 & D7, D8 \\
WP5 & Design criteria and trade-off study; \tx{DARC} specification, self-analysis, hardware benchmarking & joint & O4 & 19--48 & D9 \\
WP6 & Recruitment and training; open-source releases; reporting and dissemination & both & O5 & 1--48 & D10--D13 \\
\bottomrule
\end{tabular}
\caption{Work packages. Deliverable numbers refer to Section~5.}
\label{tab:wp}
\end{table}

\paragraph{Integration across the two sites.}
IIT~Hyderabad leads the solver-based and ML-based cryptanalysis; IIT~Bhilai leads the physical experimentation and hardware implementation.
The two groups already collaborate, with four joint publications since 2023, including hardware-validated fault attacks~\owncite{DBLP:conf/acns/KunduGSR23, Kundu_Ghosh_Aikata_Saha_2025}.
The project will run a shared code repository, a fortnightly joint meeting, and two in-person exchange visits per year.
The postdoctoral researchers at the two sites are the day-to-day integration points.

\paragraph{Risks and mitigation.}
(i) \emph{No improved attack is found on a target cipher.}
A rigorous, reproducible negative result (a proven bound from the solver models) is still a publishable security evaluation and feeds WP5 directly.
(ii) \emph{ML distinguishers do not outperform classical ones on the targets.}
The hybrid framework is designed so that the classical component stands alone; the ML part is an add-on, not a dependency.
(iii) \emph{EMFI campaigns on unrolled targets prove harder than modelled.}
The bench already exists and clock-glitch and voltage-glitch fall-backs are available; the fault-friendly distinguisher tool (D4) is independent of the physical result.
(iv) \emph{Recruitment delays.}
Both groups have a pipeline of MTech and PhD students in the area; the RA positions can start with internal candidates.
(v) \emph{Scope of the design thrust.}
\tx{DARC} starts only in month~19, with a defined fall-back of publishing the design criteria and a candidate specification if full third-party-ready documentation is not reached by month~48.

% ─────────────────────────────────────────────────────────────────────────────
\section{Expected Knowledge Outcomes, Utility, Impact \& Future Translation}
% ─────────────────────────────────────────────────────────────────────────────
The proposed project aims to advance the scientific understanding of low-latency symmetric-key cryptography through three complementary directions: automated cryptanalysis of existing designs, physical security analysis of fully unrolled implementations, and microarchitectural security of CPU.
The expected outcomes include new cryptanalytic results, reusable analysis frameworks, experimentally validated attack methodologies, and evidence-based recommendations for improving security.
Beyond individual research contributions, the project will establish methodologies and resources that can support the systematic evaluation and future development of low-latency cryptographic primitives.

\subsection{Advancement of Fundamental Knowledge}
The project will address the identified knowledge gaps through the following research outcomes, organised by research thrust and are intended to be concrete, verifiable, and directly useful to the cryptographic community and standardisation bodies.

\subsection*{Thrust 1 — Cryptanalysis of Low-Latency Ciphers}
The project will advance the understanding of the security of existing low-latency block ciphers, including \tx{ARADI}~\cite{cryptoeprint:2024/1240}, \tx{ChiLow}~\cite{DBLP:conf/eurocrypt/BelkheyarDGLMPRSSAV25}, \tx{SCARF}~\cite{DBLP:conf/uss/CanaleGLTTU23}, and \tx{QARMAv2}~\cite{DBLP:journals/tosc/AvanziBDEGNR23}.
Although automated cryptanalysis has proved effective for conventional symmetric-key primitives, its systematic application to low-latency designs, particularly for jointly optimising distinguishers and key-recovery attacks, remains an important research direction. The project will investigate this problem using constraint programming (CP), SAT, and mixed integer linear programming (MILP), alongside hybrid approaches combining machine-learning-based distinguishers with classical cryptanalysis. The expected outcomes are:



\begin{enumerate}[label=\textbf{D\arabic*.}, leftmargin=*, itemsep=2pt]
  \item \textbf{Cryptanalytic results on deployed ciphers.}
New distinguishers and key-recovery results targeting at least two of the selected ciphers, with publication targets including CRYPTO, EUROCRYPT, ASIACRYPT, or IACR Transactions on Symmetric Cryptology (ToSC).

  \item \textbf{Automated constraint-model library.}
        An open-source library of CP/SAT/MILP constraint models for the full family of low-latency block ciphers studied, enabling reproducible and extensible cryptanalysis by the wider community.

  \item \textbf{ML--classical hybrid distinguisher framework.}
        A principled methodology and associated toolchain for combining neural differential distinguishers with classical meet-in-the-middle and boomerang techniques, validated on at least one target cipher.
\end{enumerate}

\subsection*{Thrust 2 — Physical Attacks on Low-Latency Implementations}
Fully unrolled implementations execute all rounds through simultaneously instantiated hardware logic.
Their physical security characteristics differ from those of iterative implementations, particularly in terms of fault injection, fault propagation, and the spatial distribution of vulnerable components.
The project will investigate these characteristics by combining automated distinguisher searches with experimental fault analysis. The expected outcomes are:

\begin{enumerate}[label=\textbf{D\arabic*.}, leftmargin=*, itemsep=2pt, start=4]
  \item \textbf{Fault-attack-friendly distinguisher search tool.}
        A solver-based automated tool for finding fault-invariant and fault-propagation distinguishers in fully-unrolled hardware implementations, released publicly.

  \item \textbf{Experimental fault-attack results.}
        End-to-end hardware validation of theoretical fault attacks on at least one low-latency cipher, leveraging the team's existing EMFI and ChipWhisperer infrastructure, targeting publication at IACR TCHES or USENIX Security.

  \item \textbf{Hardening guidelines.}
        A set of actionable design recommendations for protecting fully-unrolled implementations against fault injection, informed by the vulnerabilities uncovered during the project.
\end{enumerate}

\subsection*{Thrust 3 — Microarchitectural Security}
The project will investigate the cryptographic components underlying CPU security mechanisms, focusing on at least one memory-encryption or pointer-authentication mechanism.
The objective is to understand how the cryptographic security of these components interacts with the microarchitectural context in which they are deployed and to identify opportunities for improvement. The expected outcomes are:
\begin{enumerate}[label=\textbf{D\arabic*.}, leftmargin=*, itemsep=2pt, start=7]
  \item \textbf{Vulnerability analysis of CPU cryptographic features.}
        A systematic security analysis of the cryptographic components underpinning at least one hardware memory-encryption or pointer-authentication mechanism, targeting top security venues (IEEE S\&P, ACM CCS, USENIX Security).

  \item \textbf{Cryptographic solutions for microarchitectural threats.}
        Concrete primitive proposals or hardening constructions for identified vulnerabilities, with security proofs or cryptanalytic evidence of their adequacy.
\end{enumerate}

These outcomes will bridge cryptographic analysis and systems-security evaluation by examining security assumptions in the context of actual hardware-assisted mechanisms.
Where weaknesses are identified, the project will investigate appropriate cryptographic improvements or hardening approaches.


\subsection{Improvement over the Scientific and Technological Baseline and Potential Utility}

The potential utility of these outcomes includes:
\begin{itemize}
    \item Supporting cryptographic researchers and designers in evaluating the security of existing and future low-latency primitives.

    \item Providing reusable, open-source tools and experimental resources that facilitate independent verification and extension of the project's findings.

    \item Helping hardware-security researchers identify vulnerable components and assess the effectiveness of implementation-level countermeasures.

    \item Informing the evaluation and potential strengthening of memory-encryption, pointer-authentication, and secure-cache mechanisms.

    \item Providing evidence-based design and hardening recommendations that account for security, latency, and implementation constraints.
\end{itemize}

The project does not assume that new practical attacks will necessarily be found against every target or that every identified weakness will lead to an exploitable vulnerability.
Improvements over the baseline may also take the form of stronger security evaluations, improved search procedures, reproducible negative results, or better-supported design recommendations.

\subsection{Long-Term Impact and Future Translation}

The project's findings may contribute to the longer-term development of more secure hardware-assisted cryptographic mechanisms.
Its potential technological impact includes improved methods for evaluating the security of primitives used in memory encryption, pointer authentication, and
cache randomisation, as well as better understanding of the trade-offs between low latency and implementation-level protection. The principal translational pathways will be:

\begin{enumerate}[label=\textbf{D\arabic*.}, leftmargin=*, itemsep=2pt, start=9]
  \item \textbf{Human resource development.}
        Training of two PhD students and two postdoctoral researchers at the intersection of symmetric-key cryptography, hardware security, and automated analysis.
        The IIT Hyderabad PhD student will work on automated and ML-based cryptanalysis, a relatively new and rapidly evolving area; the IIT Bhilai PhD student on fault analysis of unrolled hardware.
        One postdoctoral researcher at each institute will lead, respectively, the solver and GPU toolchain and the hardware experimentation and \tx{DARC} implementation.

  \item \textbf{Open dissemination.}
        All cryptanalytic models, attack tools, and experimental datasets produced during the project will be released as open-source artefacts accompanying the corresponding publications.

    \item \textbf{Design and implementation guidance:}
    The findings from the three research thrusts will be
    consolidated into recommendations for evaluating and
    hardening low-latency cryptographic designs and their
    hardware implementations.

    \item \textbf{Engagement with the cryptographic community:}
    Research publications and dissemination of reproducible
    results will support independent scrutiny and may inform
    future cipher-design revisions and standardisation
    discussions.
\end{enumerate}

In the longer term, the methodologies developed in this project may be extended to other symmetric-key constructions and security-sensitive computing environments. 
The open-source tools, validated analyses, and trained researchers will provide
a foundation for continued research and potential collaboration with cryptographic designers, hardware-security teams, and standardisation communities. 
Any direct industrial adoption will depend on subsequent evaluation and integration by the relevant stakeholders.

% \subsection*{Training and Dissemination}
% \begin{enumerate}[label=\textbf{D\arabic*.}, leftmargin=*, itemsep=2pt, start=9]
%   \item \textbf{Human resource development.}
%         Training of at least two PhD students and one postdoctoral researcher in the intersection of symmetric-key cryptography, hardware security, and automated analysis.

%   \item \textbf{Open dissemination.}
%         All cryptanalytic models, attack tools, and experimental datasets produced during the project will be released as open-source artefacts accompanying the corresponding publications.
% \end{enumerate}

% ─────────────────────────────────────────────────────────────────────────────

\section{Timeline \& Milestones}

% \section{Team, Expertise \& Institutional Facilities}
% \medskip
% \noindent\textbf{Team expertise and existing infrastructure.}
% The team brings extensive and demonstrated expertise to this thrust, spanning both the theoretical modelling of fault attacks and their physical realisation in laboratory settings.
% Prior publications on fault analysis of GIFT-family ciphers~\owncite{DBLP:conf/acns/KunduGSR23, Kundu_Ghosh_Aikata_Saha_2025} have involved end-to-end experimental validation, from distinguisher derivation through to actual key recovery on hardware.
% Crucially, an operational experimental infrastructure, including ChipWhisperer-based clock-glitch and power-analysis platforms as well as electromagnetic fault injection (EMFI) equipment, is already in place within the team.
% This means that hardware validation of theoretical fault-attack results can proceed immediately, without any capital expenditure from this proposal.
% The funding requested under this thrust is therefore directed entirely toward personnel and computational resources, rendering it a low-overhead, high-impact research direction with a clear and near-term path to results.


% \section{Justification for proposed budget}

% \section{Current Financial Support (Other Than This Proposal)}

% \section{Name and address of experts / institution interested in the subject/outcome of the project}
\section*{Conclusion}

This proposal focuses on the design and analysis of symmetric-key primitives in one critical and emerging domain: low-latency cryptography.
The investigation spans theoretical cryptanalysis, resistance to physical attacks, and the formalization of design criteria, thereby advancing the foundations of symmetric-key cryptography.
Furthermore, the fundamental expertise in symmetric-key design and cryptanalysis at the core of this proposal is directly applicable to other domains, including the development of high-performance ciphers for beyond-5G communications.

\pagebreak
\bibliographystyle{alpha}
\bibliography{../bibliography.bib,../myPublications.bib}
\end{document}
